% !PRINTED 494
\weberpage{494}

For if
% [sic] The print omits the comma between $a_1$ and $a_2$ in the lower row of (1); preserved source-faithfully.
\weq{1}{\pi = \begin{pmatrix} 0, & 1, & 2 \dots m-1 \\ a_0, & a_1 & a_2 \dots a_{m-1} \end{pmatrix},}
one begins with an arbitrary figure, say with 0, and continues the sequence
\weq{2}{0, \quad a_0 = b, \quad a_b = c \dots}
until one strikes upon some figure for the second time. The first figure to recur must be 0, since for every figure the one preceding it in the sequence (2) is uniquely determined by (1). The figures
\weq{3}{(0, b, c \dots)}
then form a first cycle. If not all $m$ figures of (1) are thereby exhausted, one picks out one of the remaining figures and proceeds in the same way, until all $m$ figures of (1) have been accommodated in the cycles. Since in every such cycle, for each figure, the preceding one as well as the following one is completely determined by (1), the cycles themselves are also uniquely determined. In the \gesperrt{designation} of $\pi$ by the cycles, however, not only may the various cycles be arranged arbitrarily, but in each cycle one may also begin with any one of its figures. A figure that is not changed forms by itself a one-term cycle.

We choose an entirely arbitrary example, by which the procedure at once becomes clear:
\[
\pi_1 = \begin{pmatrix} 0, & 1, & 2, & 3, & 4, & 5, & 6, & 7 \\ 2, & 3, & 6, & 7, & 1, & 4, & 5, & 0 \end{pmatrix} = (0, 2, 6, 5, 4, 1, 3, 7)
\]
\[
\pi_2 = \begin{pmatrix} 0, & 1, & 2, & 3, & 4, & 5, & 6, & 7 \\ 4, & 2, & 1, & 3, & 5, & 7, & 0, & 6 \end{pmatrix} = (0, 4, 5, 7, 6)(1, 2)(3)
\]
\[
\pi_3 = \begin{pmatrix} 0, & 1, & 2, & 3, & 4, & 5, & 6, & 7 \\ 7, & 2, & 3, & 6, & 5, & 0, & 1, & 4 \end{pmatrix} = (0, 7, 4, 5)(1, 2, 3, 6).
\]
The setting of the cycles side by side is equivalent to a composition in the sense used hitherto; only it is to be remarked that permutations which contain no common figure whatever can always be interchanged in the composition.

One-term cycles, which change nothing, are also often omitted in the designation.
